% Example "lorenz": the problem, the attractor and the divergence of the
% numerical solutions, and the convergence and predictability tests. Every
% number comes from data/, written by examples/lorenz/main.f90: run
% `fpm run --example lorenz` from the project root, then build this file with
% lualatex in examples/lorenz/.
\documentclass[11pt,a4paper]{article}
\usepackage[margin=2.2cm]{geometry}
\usepackage{amsmath,booktabs,array}
\usepackage{pgfplots}
\usepgfplotslibrary{groupplots,colormaps}
\pgfplotsset{compat=1.18}
\pagestyle{empty}

\colorlet{m-euler}{black}
\colorlet{m-order2}{orange!85!black}
\colorlet{m-order3}{green!50!black}
\colorlet{m-rk4}{blue!80!black}
\colorlet{m-dp54}{red!75!black}

\begin{document}
\section*{Example: the Lorenz system}

\subsection*{1. Problem}
Lorenz's model of convection (1963), listed in the SciML problem library as
\texttt{prob\_ode\_lorenz}, with his classic parameters:
\begin{gather*}
  x' = \sigma (y - x), \qquad y' = x(\rho - z) - y, \qquad z' = xy - \beta z,\\
  (\sigma, \rho, \beta) = (10, 28, \tfrac83), \qquad (x, y, z)(0) = (1, 0, 0).
\end{gather*}
The solutions settle on a strange attractor and are chaotic: two solutions
that start close separate on average like $e^{\lambda t}$, with
$\lambda \approx 0.906$ the largest Lyapunov exponent. A numerical solution
is such a nearby solution, so it follows the true one only for a limited
time, however good the method. The reference is a \texttt{dp54} run at a tiny
step, itself trusted only up to a time measured in the run.
\begin{center}
\small
% Written by examples/lorenz/main.f90: do not edit.
\begin{tabular}{l>{\raggedright\arraybackslash}p{12.4cm}}
\toprule
parameters & $(\sigma, \rho, \beta) = (10, 28, 8/3)$ \\
initial condition & $(x, y, z) = (1, 0, 0)$ \\
reference & \texttt{dp54}, $\Delta t = 9.766\times10^{-6}$, $0 \le t \le 40$; trusted to $t = 34.8$ (departs by $10^{-6}$ from the run at twice the step, plus the time to grow to $10^{-2}$) \\
convergence & error $\max_i|u_i - u_i^{\mathrm{ref}}|$ at $t = 1$, $\Delta t = 0.02, \dots, 0.02/128$; order from the two finest levels above $10^{-11}$, kept if within 0.3 of the next slope \\
predictability & $T_{\mathrm{div}}$: first $t$ with $|u - u^{\mathrm{ref}}| > 1$; slope: least squares of $T_{\mathrm{div}}$ against $\ln(1/\Delta t)$ over the levels with $1 < T_{\mathrm{div}} < t_{\mathrm{trust}} - 3$ (at least 3) \\
Lyapunov exponent & Benettin's method: two \texttt{dp54} runs ($\Delta t = 10^{-3}$), separation $10^{-8}$ restored every time unit, mean growth rate over 500 units: $\lambda = 0.901$ \\
figures & attractor: \texttt{dp54}, $0 \le t \le 60$; solutions at $\Delta t = 0.01$ \\
methods & \texttt{rk3ws}, \texttt{rk4}, \texttt{dp54} \\
interface & plain arrays, \texttt{rk\_array\_scheme\_t} \\
\bottomrule
\end{tabular}

\end{center}

\newpage
\subsection*{2. Solutions}
\begin{center}
\begin{tikzpicture}
  % ---- (a) the attractor, the classic x-z view ----
  \begin{axis}[
      name=attractor,                                      % anchor for the panels on the right
      width=8.6cm, height=11.4cm,                          % tall: z spans 0 ... 50
      xmin=-21, xmax=21, ymin=-1, ymax=52,                 % the butterfly
      xlabel={$x$}, ylabel={$z$},                          % projection on the x-z plane
      title={(a) the attractor, $0 \le t \le 60$},         % panel title
      colormap/viridis, colorbar horizontal,               % colour = time
      colorbar style={at={(0.5,-0.11)}, anchor=north, xlabel={$t$}, height=0.2cm, width=6cm},
      axis background/.style={fill=black!2},               % very light background
    ]
    \addplot[mesh, point meta=\thisrow{t}, line width=0.3pt]
      table[x=x, y=z] {data/attractor.dat};                % the trajectory, coloured by time
    \addplot[only marks, mark=*, mark size=2pt, red!80!black] coordinates {(1,0)};             % start (1, 0, 0)
    \addplot[only marks, mark=+, mark size=4pt, very thick, black] coordinates {(8.485,27) (-8.485,27)}; % equilibria
  \end{axis}
  % ---- (b) x(t): reference and two methods ----
  \begin{axis}[
      name=series,                                         % anchor for (c)
      at={($(attractor.north east)+(1.9cm,0cm)$)}, anchor=north west, % right of (a)
      width=6.4cm, height=5.4cm,                           % panel
      xmin=0, xmax=30, ymin=-22, ymax=22,                  % 30 time units
      ylabel={$x$},                                        % first variable
      title={(b) $x(t)$, $\Delta t = 0.01$},               % panel title
      axis background/.style={fill=black!2},               % very light background
      legend style={font=\tiny, at={(0.02,0.02)}, anchor=south west, fill=white, fill opacity=0.85, text opacity=1, draw=black!30},
    ]
    \addplot[black!25, line width=2.6pt] table[x=t, y=x] {data/reference.dat};            % reference
    \addlegendentry{reference}
    \addplot[m-rk4, thin] table[x=t, y=x] {data/solution_rk4.dat};                        % RK4
    \addlegendentry{\texttt{rk4}}
  \end{axis}
  % ---- (c) error growth ----
  \begin{axis}[
      at={($(series.south west)+(0,-1.6cm)$)}, anchor=north west, % below (b), room for its title
      width=6.4cm, height=6.2cm,                           % panel
      xmin=0, xmax=40, ymode=log, ymin=1e-14, ymax=1e3,    % error grows from round-off to the attractor size
      xlabel={$t$}, ylabel={$|u - u^{\mathrm{ref}}|$},     % distance to the reference
      title={(c) error growth, $\Delta t = 0.01$},         % panel title
      axis background/.style={fill=black!2},               % very light background
      legend style={font=\tiny, at={(0.5,-0.25)}, anchor=north, legend columns=2, draw=none}, % below the panel
    ]
    \addplot[m-rk4] table[x=t, y=error] {data/solution_rk4.dat};
    \addlegendentry{\texttt{rk4}}
    \addplot[m-dp54] table[x=t, y=error] {data/solution_dp54.dat};
    \addlegendentry{\texttt{dp54}}
    \addplot[black!60, dashed] table[x=t, y=distance] {data/twins.dat};                   % two exact-ish runs 1e-10 apart
    \addlegendentry{twins, $10^{-10}$ apart}
    \addplot[black, thin, densely dotted, domain=0:40] {1e-10*exp(0.906*x)};              % e^{lambda t}
    \addlegendentry{$10^{-10}e^{0.906\,t}$}
  \end{axis}
\end{tikzpicture}
\end{center}
(a) The butterfly, seen along $y$: the solution winds around the two unstable
equilibria at $x = y = \pm\sqrt{72}$, $z = 27$ (crosses), switching wing at
irregular times; colour is time. From $(1, 0, 0)$ (red) it first spirals out
around one equilibrium for about 17 time units. (b) At $\Delta t = 0.01$ \texttt{rk4} leaves the
reference near $t \approx 18$; afterwards it is still on the attractor, just
at a different place. (c) The errors grow on average like $e^{0.906\,t}$
(dotted), in bursts, until they reach the size of the attractor; the error
level the method starts from fixes how long it lasts. Growth rates over a
finite time vary along the orbit: two runs started $10^{-10}$ apart (dashed)
barely separate during the initial spiral and grow fast only after the wing
switches begin.

\newpage
\subsection*{3. Tests: convergence and predictability}

\paragraph{Prediction.} (i) Over a short time each method converges at its
order. (ii) An error $C\Delta t^p e^{\lambda t}$ reaches 1 at
$T_{\mathrm{div}} = \bigl(\ln(1/C) + p\ln(1/\Delta t)\bigr)/\lambda$: halving
the step buys $\ln 2 \cdot p/\lambda$ time units, so $T_{\mathrm{div}}$ grows
linearly in $\ln(1/\Delta t)$ with slope $p/\lambda$, and each method gives
an estimate $\lambda = p/\text{slope}$. The order $p$ here is the one that
governs how errors accumulate. For \texttt{rk3ws} that is 3, not 2: its only
third-order defect is the local error $-\tfrac{h^3}{24}f''(f,f)$, and since
$f''(f,f) = \tfrac{d}{dt}(f'f) - f'(f'f)$ along any solution of an autonomous
system, its global error is $h^2\bigl(X(u(t)) - \Phi(t,0)X(u_0)\bigr) +
O(h^3)$ with $X = -\tfrac1{24}f'f = -\tfrac1{24}u''$: a second-order term
that does not accumulate (\texttt{rk3ws} has \emph{effective order} 3).

\paragraph{Result.}
\begin{center}
\small
% Written by examples/lorenz/main.f90: do not edit.
\begin{tabular}{lcccccc}
\toprule
name & order & observed & $T_{\mathrm{div}}$ at $\Delta t = 0.01$ & slope of $T_{\mathrm{div}}$ & $p_{\mathrm{eff}}$ & $\lambda = p_{\mathrm{eff}}/\mathrm{slope}$ \\
\midrule
\texttt{rk3ws} & 2 & 2.25 & 14.77 & 3.40 & 3 & 0.88 \\
\texttt{rk4} & 4 & 3.92 & 18.11 & 5.26 & 4 & 0.76 \\
\texttt{dp54} & 5 & 4.96 & 24.50 & 5.82 & 5 & 0.86 \\
\bottomrule
\end{tabular}

\end{center}

\begin{center}
\begin{tikzpicture}
  \begin{groupplot}[
      group style={group size=2 by 1, horizontal sep=1.8cm}, % two panels
      width=7.4cm, height=6.6cm,                           % size of each panel
      grid=major, grid style={black!8},                    % faint grid
      axis background/.style={fill=black!2},               % very light background
      every axis plot/.append style={thick, mark size=1.3pt}, % same weight for all curves
    ]
    % ---- (i) convergence ----
    \nextgroupplot[
      xmode=log, ymode=log, title={(i) error at $t = 1$}, xlabel={$\Delta t$},
      legend to name=lzlegend, legend columns=7,           % legend below the panels
      legend style={font=\small, /tikz/every even column/.append style={column sep=0.25cm}},
    ]
    \addplot[m-order3, dashed, mark=triangle] table[x=dt, y=error] {data/error_rk3ws.dat};
    \addlegendentry{\texttt{rk3ws}}
    \addplot[m-rk4, mark=diamond*]            table[x=dt, y=error] {data/error_rk4.dat};
    \addlegendentry{\texttt{rk4}}
    \addplot[m-dp54, mark=pentagon*]          table[x=dt, y=error] {data/error_dp54.dat};
    \addlegendentry{\texttt{dp54}}
    % ---- (ii) predictability ----
    \nextgroupplot[
      title={(ii) $T_{\mathrm{div}}$ against $\ln(1/\Delta t)$}, xlabel={$\ln(1/\Delta t)$},
      ylabel={$T_{\mathrm{div}}$}, ymin=0, ymax=36,        % divergence times
    ]
    \addplot[m-order3, dashed, mark=triangle] table[x=logdt, y=tdiv] {data/divergence_rk3ws.dat};
    \addplot[m-rk4, mark=diamond*]            table[x=logdt, y=tdiv] {data/divergence_rk4.dat};
    \addplot[m-dp54, mark=pentagon*]          table[x=logdt, y=tdiv] {data/divergence_dp54.dat};
    \addplot[m-rk4, thin, densely dotted]     table[x=logdt, y=tdiv] {data/divtheory_rk4.dat};    % slope 4/lambda
    \addplot[m-dp54, thin, densely dotted]    table[x=logdt, y=tdiv] {data/divtheory_dp54.dat};   % slope 5/lambda
    \addplot[black!40, thin, dashed, domain=3.5:9.2] {34.8};                                     % reference trusted below
    \node[font=\tiny, anchor=south west, text=black!60] at (axis cs:3.6,34.8) {reference trusted below};
  \end{groupplot}
  \node[below=0.2cm] at (current bounding box.south)        % legend centred under the panels
    {\pgfplotslegendfromname{lzlegend}};
\end{tikzpicture}
\end{center}
(i) Over one time unit every method converges at about its order
(\texttt{rk3ws}: 2.25, still on its way to 2). (ii) $T_{\mathrm{div}}$ climbs
in steps rather than along a line: the error grows in bursts at the wing
switches, so a run departs at one of a few critical moments, and halving the
step often only moves the departure to the same moment or the next one. On
average the climb follows the predicted slope $p/\lambda$ (dotted lines,
$\lambda$ from Benettin's method): the estimates of $\lambda$ lie between
$0.76$ and $0.88$ for the three methods, against $0.909$ from Benettin's
method. \texttt{rk3ws} yields its estimate only with $p = 3$, its effective
order.
\end{document}
